\documentclass[11pt]{article}
\usepackage[margin=0.87in]{geometry}
\usepackage{amsmath,amssymb,amsthm,booktabs,array,microtype,xcolor,hyperref}
\definecolor{accent}{HTML}{174A53}
\hypersetup{colorlinks=true,linkcolor=accent,urlcolor=accent,citecolor=accent,pdftitle={The exponent-seven census for imaginary quadratic class groups}}
\newtheorem{theorem}{Theorem}
\newtheorem{lemma}{Lemma}
\newcommand{\Cl}{\operatorname{Cl}}
\newcommand{\Q}{\mathbb Q}
\newcommand{\Z}{\mathbb Z}
\newcommand{\p}{\mathfrak p}
\newcommand{\base}{X_0}
\setlength{\parindent}{0pt}
\setlength{\parskip}{5pt plus 1pt}
\setlength{\emergencystretch}{2em}
\renewcommand{\labelitemi}{\(\bullet\)}
\newcommand{\RatioPairs}{1,963,144}\newcommand{\RatioCandidates}{231,471}\newcommand{\RatioWitness}{97}
\begin{document}
{\small\color{accent}\textsc{Research note}\hfill November 25, 2026}\par
{\LARGE\bfseries The exponent-seven census for imaginary\\[2pt] quadratic class groups}\par
{\large Split primes, seventh-power norm equations and Dirichlet GRH}\par\medskip
\begin{theorem}\label{thm}
Exactly 33 imaginary quadratic fields have ideal class group of exponent seven.
For 31 the group is cyclic of order seven. For discriminants $-63499$ and
$-118843$ it is $(\Z/7\Z)^2$. Exactly 40 imaginary quadratic orders have proper
ideal class group of exponent seven, of which seven are nonmaximal.
\end{theorem}
The members of both lists were known. Elsenhans, Klüners and Nicolae (EKN)
proved the field list complete through $\base=3.1\cdot10^{20}$ and, under ERH,
bounded every remaining discriminant by $3.9\cdot10^{30}$ \cite{EKN}.
Aiken, Ross, Shvydkoy and Xue singled out this same interval in their
order-classification theorem \cite{ARSX}. An incoming ProveIt report lowered
the upper limit to $4\cdot4600^7$ and asked whether pairs of seventh-power norm
equations could close the interval \cite{prior}. We eliminate it.
The new ingredients are the use of reduced ratios of norm coefficients and
an upper bound that charges each allowed group of coefficients only once.

\paragraph{What comes from the Lean proof.}
In the readable edition at commit \texttt{72fead10}, \texttt{Statements.lean}
defines the Dirichlet zero assertions in Mathlib terms. The lattice comparison
in \texttt{MainInequality.lean} gives \texttt{star\_inequality}.
\texttt{OneSidedBound.lean} converts it into a one-sided estimate on a smoothed
prime sum, and \texttt{Zeros/Criterion.lean} detects zeros through the residues of
its transform. In \texttt{Main.lean}, \texttt{LFunction\_eq\_zero\_re}
places all Dirichlet zeros in the open strip on the critical line.
Sutherland records independent kernel and statement checks for the pinned
source \cite{lean,verify}. We use the quadratic-character instance to give
a definite sign to the zero side of the explicit formula. The norm equations
then constrain the prime side in a way that zero location alone does not supply.
All analytic hypotheses used in the census concern $\zeta$ and quadratic
Dirichlet $L$-functions. The new field classification is a mathematical proof
with finite computer checks. There is no Lean formalization of the new steps
in this archive.

\paragraph{Outline of the elimination.}
If $x=|D|>\base$, a small split prime $p$ supplies an integer $b_p\ge1$
with $x b_p^2\le4p^7$. Certain pairs of these integers cannot occur together.
An explicit formula gives a lower bound on the total contribution that split
primes have to make. The pair exclusions give a smaller upper bound. Both
inequalities are strict in the required direction even at the lower end of
the interval.

\newpage
\section{The arithmetic restrictions on split primes}
Let $D=-x$ be a fundamental discriminant with $x>\base$ and suppose
$\Cl(\Q(\sqrt D))$ has exponent seven. This finite abelian group is a
vector space over $\mathbb F_7$, so its order is odd. Genus theory gives
$2$-rank $t-1$ for $t$ distinct prime discriminant factors. Thus there is
just one ramified prime. The dyadic fundamental discriminants with this
property are small, so $x$ is an odd prime with $x\equiv3\pmod4$.
In particular all primes in our test supports $p\le7000$ are unramified.

At a split prime $p$, write $(p)=\p\bar\p$. A generator of the integral
ideal $\p^7$ has norm $p^7$ and is $(a+b\sqrt D)/2$ with $a,b\in\Z$ of
the same parity. Since its valuations at $\p$ and $\bar\p$ differ, $b\ne0$.
Taking absolute values of the coefficients gives
\begin{equation}\label{norm}
 a_p^2+x b_p^2=4p^7,\qquad a_p\ge0,\quad b_p\ge1.
\end{equation}
In particular every prime $r$ with $4r^7\le\base$ is inert. This includes
$2$, so $D\equiv5\pmod8$ and $x\equiv3\pmod8$.

\begin{lemma}[Finite exclusions]\label{finite}
For split primes $p,q$ of a putative field as above, the following hold.
\begin{enumerate}
\item Distinct $p,q\le6500$ cannot have $b_p=b_q$.
\item For distinct $p,q\le4200$, the fraction $b_p/b_q$ in lowest terms
cannot have both numerator and denominator at most 12.
\end{enumerate}
\end{lemma}
\begin{proof}
For the first assertion take $p<q$. The norm equations show that $a_p$
and $a_q$ have the same parity. Put
$u=(a_q-a_p)/2$, $v=(a_q+a_p)/2$. Then $uv=q^7-p^7$ and $0<u\le v$.
Enumerate the divisors $u\le\sqrt{q^7-p^7}$, and put
$M=4p^7-(v-u)^2$. A candidate has $M=xb_p^2>\base$.
If $-M$ is a nonzero square modulo any odd prime $r$ with
$4r^7<\base$, then $r\nmid b_p$ and $(D/r)=1$, contradicting inertness.
For each eligible pair we factor $(q-p)$ and $(q^7-p^7)/(q-p)$, enumerate
all divisors of their product, and perform this test. There are 717 eligible
primes, 256,686 pairs and 191,603 divisor candidates above the threshold.
Every candidate has a witness $r\le107$.

For unequal coefficients write $b_p=ic$, $b_q=jc$, with
$1\le i<j\le12$ coprime. Both assignments of the distinct primes occur in
the enumeration. Put $d=i^2q^7-j^2p^7$. It is nonzero, since $p,q>693$.
The two nonnegative numbers $i a_q,j a_p$ have squares differing by $4d$,
so their sum and difference are even. For every possible solution the half
sum $v$ and half absolute difference $u$ satisfy $uv=|d|$.
If $d>0$ then $j a_p=v-u$ and $i a_q=v+u$. If $d<0$ the two expressions
are exchanged. Test divisibility by $j$ and $i$, recover the two integers,
and keep only common integral values
\[
        M=(4p^7-a_p^2)/i^2=(4q^7-a_q^2)/j^2>\base .
\]
The same residue test applies because $M=xc^2$. The finite check on the
next page completes the proof.
\end{proof}

\newpage
\section{Checking the pairs of norm equations}
The program \texttt{checks/ratio\_norms.py} first discards cases in which
$4p^7\le\base i^2$ or $4q^7\le\base j^2$. It then checks the two norm
equations simultaneously modulo
\[
             128,\ 81,\ 125,\ 49,\ 121,\ 13,\ 17,\ 19,\ 23 .
\]
For an odd base prime $r$ choose a unit $w$ with $(-w/r)=-1$.
From $(-x/r)=-1$ we get that $x/w$ is a square modulo $r$. A square root lifts to
each $r^k$ by Hensel's lemma. At $2$ use $w=3$, since a dyadic unit
is a square precisely when it is $1\bmod 8$.
Thus every possible $M=xc^2$ lies in the finite set $w(\Z/r^k\Z)^2$
of square multiples, including multiples of $r$. For each modulus keep
a pair if there is a common residue $t$ in this set such that
$4p^7-i^2t$ and $4q^7-j^2t$ are both squares. This test is necessary.
The choices of $t$ may be made separately at the different moduli.

For each surviving pair, every divisor $u$ of $|d|$ up to $\sqrt{|d|}$
is tested as in Lemma~\ref{finite}. A final witness $r>12$ with $4r^7<\base$
and $-M$ a nonzero square suffices. The aggregate results of the complete
run are as follows.
\begin{center}
\begin{tabular}{lrr}\toprule
 & Equal coefficient & Distinct coefficients\\\midrule
Prime pairs processed after preliminary tests & 256,686 & \RatioPairs\\
Divisor candidates giving a possible $M>\base$ & 191,603 & \RatioCandidates\\
Largest first witness $r$ & 107 & \RatioWitness\\
Candidates left & 0 & 0\\\bottomrule
\end{tabular}
\end{center}
The first column counts unordered pairs. The second counts ordered prime
pairs for each reduced ratio, with multiplicity over ratios. Running
\texttt{equal\_norms.py} and \texttt{ratio\_norms.py} reproduces the columns.
The second program reports counts for every ratio in its output.
The divisor searches use exact integer arithmetic and the proof-enabled Debian build of GNU \texttt{factor} 9.7. The binary hash is recorded in the archive. Each list of factors is checked against its product.
NumPy packs exact residue masks and enumerates candidate pairs. No
floating-point decision enters either exclusion.

\paragraph{Why use ratios.} Suppose one coefficient equals $2c$ and
another equals $3c$. Neither $c$ nor the discriminant need be known.
The pair would force a factorization of $|4q^7-9p^7|$, which depends only
on the primes, and one of its divisors would have to produce an admissible
$M$. Clearing both coefficients in this way makes a single calculation
rule out every value of $c$. Limiting the reduced multipliers to 12 is a
choice of finite search size. The interval inequalities below verify it suffices.

\newpage
\section{A lower bound from the zeros}
Write $\chi=(D/\cdot)$, an odd primitive quadratic character of
conductor $x$. Let $g$ be a compactly supported real step function,
$F(u)=\int g(v)g(v+u)\,dv/\int g(v)^2\,dv$, and
$h(t)=\int F(u)e^{itu}\,du$. Thus $F(0)=1$ and $h(t)\ge0$ for real
$t$. For our step functions $h$ is entire and decays quadratically on
each fixed horizontal strip. The Guinand--Weil formula \cite{explicit} gives
\[
 \sum_\rho h((\rho-1/2)/i)=\log(x/\pi)+
   \frac1{2\pi}\int h(t)\Re\psi(3/4+it/2)\,dt
       -2\sum_{n\ge2}\frac{\Lambda(n)\chi(n)}{\sqrt n}F(\log n).
\]
Here $\psi=\Gamma'/\Gamma$. The sum over nontrivial zeros is nonnegative by
Dirichlet GRH. This is the only use of zero location in the window calculation.
If $C$ is the largest absolute slope of $F$, Fourier inversion in the
digamma integral gives the archimedean term as
\[
 \psi(3/4)+2\int_0^\infty\frac{e^{-3u/2}}{1-e^{-2u}}(1-F(u))\,du
 \le\psi(3/4)+\tfrac{14}{9}C .
\]
Indeed $1-F(u)\le Cu$ for $u\ge0$, and the geometric expansion gives
$\sum_{k\ge0}(2k+3/2)^{-2}\le4/9+1/3$. Also
$\log\pi-\psi(3/4)>2$, using $\log\pi>1.1$,
$\gamma>H_6-\log7>1/2$, and
$\psi(1)-\psi(3/4)>1/3+1/14$ from the first two summands of its series.
Consequently
\begin{equation}\label{weil}
 \log x>2-14C/9+2\sum_{n\ge2}\frac{\Lambda(n)\chi(n)}{\sqrt n}F(\log n).
\end{equation}

For each $Y$ use 12 consecutive intervals of width $\log(Y)/12$ for $g$.
The last six heights are the negatives of the first six in reverse order. Use either first half
\[
 V_0=(21,14,10,6,4,1),\qquad V_1=(1000,933,686,449,251,81).
\]
$V_0$ and the archimedean estimate were used in the incoming report.
$V_1$ was chosen numerically and the resulting rational test was checked
exactly. The support of $F$ is $[-\log Y,\log Y]$. Define its inert-prime
baseline and its maximum loss at a split prime by
\begin{align*}
 S(Y)&=2-14C/9+
   2\sum_{p^k\le Y}(-1)^k\frac{\log p}{p^{k/2}}F(k\log p),\\
 W_Y(p)&=\max\{0,-4(\log p)F(\log p)/\sqrt p\}.
\end{align*}
Every prime is unramified in the supports. A split prime is greater than
693 and its square exceeds every $Y$ used here. Hence changing its sign
relative to the inert model subtracts at most $W_Y(p)$ from $S(Y)$.
If the total charge is at most $L$, equation~\eqref{weil} forces
$\log x>S(Y)-L$. A contradiction follows as soon as $S(Y)-L>\log U$
on a proposed interval $x\le U$.

\newpage
\section{Charging the split primes}
Let $A<x\le U$ and put $B_p=\lfloor\sqrt{4p^7/A}\rfloor$.
Equation~\eqref{norm} implies $b_p\le B_p$. For $Y\le6500$, distinct
split primes must use distinct integers $b_p$. At $Y\le4200$ they must
also avoid the ratios in Lemma~\ref{finite}. The following partition uses
both observations. Set $B=\max_{p\le Y} B_p$. Process $1,\ldots,B$ in
order. At the first unassigned $b$ form a block of all still unassigned
members of $\{b,2b,\ldots,kb\}$ at most $B$. Let $I$ be the set of block
minima. We use either $k=1$ or $k=12$.

Two members of the same block have ratio in lowest terms with numerator
and denominator at most $k$. At most one split prime can therefore use
any given block, including the possibility of two equal coefficients.
If its minimum is $b$, that prime must satisfy $B_p\ge b$. Hence
\begin{equation}\label{pack}
 L(A,Y,k)=\sum_{b\in I}\ \max_{p\le Y,\ B_p\ge b}W_Y(p)
 \quad\text{is an upper bound for the total split-prime charge.}
\end{equation}
Each term is a separate maximum, so allowing the same prime to realize
several maxima preserves the upper bound. This argument also permits a
block for which no split prime exists.

\paragraph{Reaching the final interval.}
Start from the EKN bound $U=3.9\cdot10^{30}$, which uses the split-prime
bound of Bach and Sorenson for a quadratic field \cite{Bach}. Its
Dedekind zeta function is $\zeta(s)L(s,\chi)$.
With heights $V_0$, test $Y=7000,5000,4700,4600$ in sequence.
If $x>4Y^7$, every prime through $Y$ is inert, hence $L=0$. Each row
forces a lower bound on $\log x$ exceeding the logarithm of its current
ceiling. Update that ceiling to $4Y^7$. Applying \eqref{pack} next
excludes the two intervals with lower endpoints $4\cdot10^{22}$ and
$10^{22}$. The rounded bounds in the table are outward.
\begin{center}
\begin{tabular}{rlrrr}\toprule
$Y$ & Lower endpoint $A$ & Charge upper bound & $S-L$ lower bound & $\log U$ upper bound\\\midrule
7000 & $4\cdot7000^7$ &0&76.072&70.439\\
5000 & $4\cdot5000^7$ &0&63.631&63.362\\
4700 & $4\cdot4700^7$ &0&61.571&61.007\\
4600 & $4\cdot4600^7$ &0&60.867&60.574\\
6500 & $4\cdot10^{22}$ &12.055&61.084&60.423\\
4900 & $10^{22}$ &10.186&52.767&52.044\\\bottomrule
\end{tabular}
\end{center}
The four inert-only rows reproduce the incoming exponent-seven bound.
The next two use equal coefficients. The eight rows on the next page
use the unequal-coefficient exclusion with $Y=4200$ and heights $V_1$.
Together the fourteen intervals cover every $x>\base$ allowed by EKN.

\newpage
\section{Closing the interval and recovering the census}
Each row below excludes $A<x\le U$. The first two columns give the
endpoints in units of $10^{20}$. The fourth column rounds down
$S(4200)-L(A,4200,12)$. The last rounds up $\log U$.
\begin{center}
\begin{tabular}{rrrrr}\toprule
$A/10^{20}$ & $U/10^{20}$ & Charge upper bound & Forced lower bound & Ceiling\\\midrule
3.1&3.255&9.223&47.295&47.232\\
3.255&3.74325&9.001&47.518&47.372\\
3.74325&4.866225&8.368&48.151&47.635\\
4.866225&7.2993375&7.322&49.197&48.040\\
7.2993375&13.1388075&6.011&50.508&48.628\\
13.1388075&26.277615&4.600&51.919&49.321\\
26.277615&78.832845&3.425&53.094&50.420\\
78.832845&100&2.120&54.399&50.657\\\bottomrule
\end{tabular}
\end{center}
At the first endpoint there are 193 blocks. Even if every block uses
the most expensive prime available to it, the total charge is below
$9.223$. The inert baseline exceeds $56.518$. Their difference exceeds
the upper bound $47.232$ for $\log x$. This is the closest comparison.
The seven subsequent comparisons have larger margins.

\paragraph{Reproducibility of the analytic check.}
At knots $u=j\log Y/12$, $F$ equals
$\sum_i v_i v_{i+j}/\sum_i v_i^2$, with the full signed vector extended
by zero. It is linear between knots. For $n=p^a$ locate the knot by
comparing $n^{12}$ to $Y^j$ as integers. To enclose $\log n$, write
$n=2^b t$, $1\le t\le2$, use $w=(t-1)/(t+1)$, and bound the tail after
the first 19 terms of $2\sum_{j\ge0}w^{2j+1}/(2j+1)$ by
$2w^{39}/[39(1-w^2)]$. The integer square root of $10^{28}n$ bounds
$1/\sqrt n$ on both sides. The script uses fractions and outward rounding
to a $10^{-14}$ grid. A separate 65-digit calculation recomputes the
baseline by overlapping translated intervals directly. The exact script
is the decision procedure for the strict inequalities.

\paragraph{Fields and orders.}
EKN's bounded census now gives the first assertion of Theorem~\ref{thm}.
For the second, apply \cite[Theorem 1.3(6), Section 6]{ARSX}.
Their conductor enumeration starts from the maximal orders whose class-group
exponent divides seven. The class-number-one fields are known, and the
field calculation above discharges the remaining completeness condition.
The endpoint $3.1\cdot10^{20}$ itself is not fundamental. The seven
additional order discriminants are
\[
 -284,\ -343,\ -604,\ -892,\ -1372,\ -1852,\ -1948.
\]
The 38 cyclic groups and two rank-two groups in Table A.7 have together
364 classes. The supplied reduced-form calculation checks the seventh power
of each class and records generators. The CSV gives the field discriminant,
conductor, class number and group for each row. This completes the proof.

\newpage
\section{Other transfers from the proof}
Four companion notes address different questions. The finite census above
uses the endpoint of the Dirichlet proof. The first three notes also examine
internal constructions that can be studied on their own.

\begin{itemize}
\item \emph{The cost of simultaneous torsion specialization over abelian local
fields.} Starting with \texttt{RootsOfUnity.lean}, the note determines the exact
valuation budget needed to obstruct every permitted primitive-root evaluation
of a family of integral polynomials. It covers all finite abelian coefficient
extensions of $\Q_p$, all finite choices of orders and each nonnegative
tolerance. Tame descent and an odd-degree approximation argument at $2$
extend the cyclotomic and fourth-root cases in ProveIt. This result uses no
hypothesis about zeros.
\item \emph{The sparse-place cutoff.} A refinement of the digit construction
in the source file \texttt{Existence.lean} gives a large set that works at
every prime for which even one place can work. For the stated instance
$q_0=18480$, $s=44$, 120 places work up to the first obstruction,
the prime $18\,216\,708\,648\,992\,767$. A finite sieve and a primality
certificate verify the cutoff.
\item \emph{Recovering the modulus and quadratic twist from one finite band
of signed zeros.} For the lattice parameters in the readable proof, a
sufficiently high finite band of the aggregated signed multiplicities
determines both integers with a stopping rule. The data include certified
error budgets. The finite example certifies sufficient precision and band
height for distinguishing two packets; it does not compute their zero lists.
\item \emph{Exceptional cutoffs for the central Euler product.} For a fixed
primitive nonprincipal Dirichlet character, the measure in $[U,U+1]$ of
logarithmic cutoffs with normalized canonical-log error exceeding $\epsilon$
is $O_{\chi,\epsilon}(e^{-\sqrt{2\pi\epsilon U}})$. This gives a quantitative
partial result for the deep Riemann hypothesis, leaving the all-cutoff
limit open.
\end{itemize}
\paragraph{Further questions.}
For exponent seven, what is the smallest maximum reduced multiplier
that suffices after optimizing the test and the partition? Can the
ratio exclusion be verified by a shorter, formally checked arithmetic
certificate? For higher exponents one can test the same general norm
equations, but a new initial census range, prime weights and compatibility
proof would be needed. The dyadic local result also suggests asking
which nonabelian coefficient fields have the same tolerance profile after
a tame base change.


\newpage
\section*{The discriminants}
The following are the absolute values of the 33 field discriminants. A
minus sign gives the actual discriminant. The 31 entries in the first row
have class number seven.
\begin{center}
\begin{tabular}{lp{0.70\linewidth}}\toprule
Class group & $|D|$ for the maximal orders\\\midrule
$\Z/7\Z$ & 71, 151, 223, 251, 463, 467, 487, 587, 811, 827, 859, 1163, 1171, 1483, 1523, 1627, 1787, 1987, 2011, 2083, 2179, 2251, 2467, 2707, 3019, 3067, 3187, 3907, 4603, 5107, 5923 \\[3pt]
$(\Z/7\Z)^2$ & $63499,\ 118843$ \\\bottomrule
\end{tabular}
\end{center}
\begin{center}
\begin{tabular}{rrr}\toprule
Nonmaximal $D$ & Field discriminant & Conductor\\\midrule
$-284$ & $-71$ &2\\
$-343$ & $-7$ &7\\
$-604$ & $-151$ &2\\
$-892$ & $-223$ &2\\
$-1372$ & $-7$ &14\\
$-1852$ & $-463$ &2\\
$-1948$ & $-487$ &2\\\bottomrule
\end{tabular}
\end{center}
In each row $D=f^2 D_K$ and the proper class group is cyclic of order
seven. A second finite computation enumerated all 59,421 negative order
discriminants of absolute value at most 118,843. Of the 135 orders with
nontrivial class number a power of seven, exactly the 40 recorded here have
exponent seven. EKN supplies the larger published census up to $\base$.
\begin{thebibliography}{9}\small
\bibitem{lean} L. Sernau and collaborators. \href{https://github.com/lsernau/grh-lean/tree/72fead102f663d5c7a407964a13dc0cf9d977207}{GRH for Dirichlet $L$-functions, in Lean 4}, readable edition, based on the proof of L. Alpöge with Claude.
\bibitem{verify} A. Sutherland. \href{https://math.mit.edu/~drew/shortgrh/}{Verification files for the shorter Lean proof of RH and GRH}, records for commit \texttt{72fead10}.
\bibitem{EKN} A.-S. Elsenhans, J. Klüners and F. Nicolae. \href{https://doi.org/10.4064/aa180220-20-3}{Imaginary quadratic number fields with class groups of small exponent}. Acta Arith. 193 (2020), 217--233. Theorems 1, 3 and Tables 1, 2. \href{https://arxiv.org/abs/1803.02056}{Preprint}.
\bibitem{ARSX} D. Aiken, E. Ross, D. Shvydkoy and H. Xue. \href{https://arxiv.org/abs/2605.26020}{Boundary CM points and class groups of small exponent} (2026), Theorem 1.3(6) and Table A.7.
\bibitem{prior} V. Reshetnikov, ProveIt collection. \href{https://github.com/VladimirReshetnikov/ProveIt/blob/d31039a48d963692e5a5176f47b6603fc9e9d9cf/docs/incoming/grh-exponent-six-and-transfers.zip}{Exponent six and transfers}, in particular \texttt{global-bound.tex} and the first question in \texttt{closing.tex}.
\bibitem{Bach} E. Bach and J. Sorenson. \href{https://doi.org/10.1090/S0025-5718-96-00763-6}{Explicit bounds for primes in residue classes}. Math. Comp. 65 (1996), 1717--1735, Table 3.
\bibitem{explicit} E. Carneiro, M. B. Milinovich, E. Quesada-Herrera and A. P. Ramos. \href{https://arxiv.org/abs/2404.08380}{Fourier optimization, the least quadratic non-residue, and the least prime in an arithmetic progression}, Lemma 9.
\end{thebibliography}
\end{document}
